\section{实现细节与实验设置}
\subsection{实现概览}
本文实现的重点集中在特征级可信建模流程。整体过程包括四个步骤：第一，利用 \encoder{} 对临床长文本进行编码，并通过均值池化获得固定维度特征；第二，在缓存特征上训练线性头、MLP 头和 Ridge 头作为比较基线；第三，训练 \method{} 的阶段 A 模型，通过 Beta-NLL 联合学习均值分支和方差分支；第四，在阶段 B 中利用方差诱导的样本权重执行 WLS 精修，对均值分支最后一层进行结构化重估。

这种实现方式把“长文本编码”和“可信输出层训练”解耦开来，从而显著降低了实验成本。编码器只需在特征提取阶段前向运行一次，之后的多组基线比较、超参数尝试和消融分析都可以在缓存特征上完成。这一策略不仅提高了训练效率，也使不同输出层方法之间的比较更具可解释性，因为所有方法都建立在同一套固定表示之上。

\subsection{数据与任务设置}
本文以 Phenotype 多标签分类任务为主实验场景。该任务使用临床文本作为输入，以十四维多热向量作为输出，适合验证长文本编码和可信输出层之间的协同关系。之所以优先选择这一任务，主要基于以下考虑：第一，多标签设置更容易体现不同标签维度上的不确定性差异；第二，临床病历中经常存在多种共病或风险因素并存的情况，任务形式更贴近真实应用；第三，该任务具有较成熟的实验组织方式，便于在统一设置下比较不同输出层。

在数据处理层面，文本和标签首先完成对齐、清洗与切分，并构造训练集、验证集和测试集。本文不再将这些预处理过程作为主体展开，而只保留与方法复现直接相关的信息：所有基线与改进方法共享同一套训练、验证和测试划分。

\subsection{特征提取设置}
特征提取阶段调用冻结的 \encoder{} 对输入文本进行编码，最终采用均值池化获得固定长度表示。之所以采用均值池化，是因为该策略实现简单、稳定性较好，并能在不增加额外参数的前提下汇聚长文本中的整体语义信息。尽管该策略并非唯一选择，但它为本文提供了一个干净且易于复现的基础表示方案。

特征提取完成后，训练集、验证集和测试集表示被缓存下来。在此基础上，后续所有输出层训练与评估都只依赖这些固定特征，而不再重复运行编码器。

\subsection{基线设置}
为评估 \method{} 的实际价值，本文设置了三类特征级基线：线性头、MLP 头与 Ridge 头。线性头用于刻画最简单的确定性映射能力；MLP 头用于检验加入非线性变换后性能是否提高；Ridge 头则用于比较带显式正则约束的结构化线性输出层。三类基线都基于相同特征和相同数据划分训练，从而保证比较公平。

在线性头和 MLP 头中，优化器均采用 AdamW，损失函数采用 BCEWithLogitsLoss；对于 Ridge 头，则在验证集上搜索正则系数，并以 MSE 为主要选择依据。通过这样的设置，本文能够比较清楚地区分“非线性提升”“结构约束提升”和“可信输出层提升”之间的差异。

\subsection{\method{} 训练设置}
在 \method{} 的训练中，阶段 A 使用 Beta-NLL 联合优化均值分支与方差分支，阶段 B 使用 WLS 对最后线性层进行精修。参考实验中，隐藏层维度设置为 256，dropout 为 0.1，学习率为 $1\times10^{-3}$，批大小为 256，Beta-NLL 中的 $\beta$ 参数取 0.5，WLS 正则系数 $\lambda$ 取 1.0。训练过程中在验证集上监控 weighted F1 与 MSE，并采用早停策略保留最优模型状态。

这一设置遵循一个基本原则：在不过度扩大模型复杂度的前提下，优先保证双分支训练的稳定性与 WLS 精修的可执行性。因此，本文并未在当前阶段引入复杂的联合微调或多级校准过程，而是将重点放在“均值分支、方差分支和结构化求解器能否形成有效配合”这一核心问题上。

\subsection{评价指标}
本文采用五类指标评价模型性能：weighted F1、micro F1、macro AUROC、MSE 和 MAE。其中，weighted F1 用于衡量整体标签分布下的分类表现，micro F1 强调所有标签实例的总体正确性，macro AUROC 则用于观察模型对不同标签维度的排序能力；MSE 与 MAE 则从连续预测分数的角度评估输出误差。采用这组指标的原因在于，\method{} 同时具有分类和不确定性建模属性，因此仅报告单一分类指标不足以反映其整体行为。

此外，为了分析可信性输出，本文还关注预测方差的均值、中位数和分布形态，并把这些统计量作为实验分析的重要辅助证据。换言之，本文的实验目标不仅是比较“谁的 F1 更高”，还包括比较“谁的输出更稳定、更可解释”。

\subsection{实现边界}
需要指出的是，本文当前实现仍然具有明确边界。首先，完整实验依赖对原始数据集的规范整理，因此跨环境复现仍需要保证数据准备过程一致。其次，本文的方法验证主要建立在已有长上下文编码器之上，关注点集中在下游输出层，而非继续预训练本身。再次，当前 \method{} 首先面向特征级多标签任务，尚未完全扩展到所有 token 级任务。明确这些边界有助于避免过度泛化，也使结论更符合真实研究状态。
